\documentclass[11pt]{article}
\usepackage[margin=1in]{geometry}
\usepackage{setspace}
\usepackage{hyperref}
\usepackage{booktabs}
\usepackage{enumitem}
\usepackage{amsmath}
\usepackage{graphicx}

\title{ORIE 5270 Final Project Summary:\\Generative Retrieval for Multi-Destination Trip Prediction}
\author{Althea Lam}
\date{\today}

\begin{document}
\maketitle
\onehalfspacing

\section{Introduction}
This project studies next-destination prediction for multi-city travel itineraries using the Booking.com Multi-Destination Trips dataset from the WSDM WebTour 2021 challenge. The practical goal is to complete incomplete trips by predicting the final city, evaluated with Accuracy@4 (whether the true next city appears in the top 4 recommendations).

Unlike standard sequence classification tasks, this problem combines three challenges:
\begin{itemize}[leftmargin=1.5em]
    \item \textbf{Large candidate space}: thousands of candidate city IDs.
    \item \textbf{Behavioral variability}: users have diverse route styles, lengths, and cross-border patterns.
    \item \textbf{Prefix-only prediction}: the model must infer the next city from partial trajectories.
\end{itemize}

The project explores three main pipelines: (1) direct city-token embedding prediction, (2) RQ-VAE code prediction with city decoding, and (3) residual k-means (RQKMeans) code prediction with decoding.

\section{Dataset and Task Setup}
The training split contains approximately 1.17 million booking rows and the test split contains 378k rows, corresponding to about 70k trips to complete. Each row includes fields such as \texttt{city\_id}, \texttt{checkin}, \texttt{checkout}, \texttt{booker\_country}, \texttt{device\_class}, \texttt{affiliate\_id}, and \texttt{hotel\_country}.

\subsection{Trip Construction}
Raw bookings are grouped by \texttt{utrip\_id} after chronological sorting. For each trip, we construct:
\begin{itemize}[leftmargin=1.5em]
    \item City sequence (\texttt{city\_id} list) as the core behavioral signal.
    \item Temporal features: month and stay-duration statistics.
    \item User/context features: booker country, device class, affiliate source.
    \item Spatial/cross-border features from hotel-country transitions.
\end{itemize}

\subsection{Training Objective}
Given a prefix of visited cities, predict the next city (or code pair) and decode to top-4 city candidates. The project supports both:
\begin{itemize}[leftmargin=1.5em]
    \item \textbf{Single-step training}: one training sample per trip.
    \item \textbf{Multi-step training}: multiple prefix-to-next samples per trip.
\end{itemize}
Empirically, multi-step training is usually beneficial, but it significantly increases sample count and runtime.

\section{Modeling Pipelines}
\subsection{Pipeline A: Direct Embedding Prediction}
The strongest baseline predicts city IDs directly from prefix tokens. Two sequence backbones were tested:
\begin{itemize}[leftmargin=1.5em]
    \item Transformer encoder variants (pooling: \texttt{last}, \texttt{mean}, \texttt{cls}).
    \item GRU next-city classifier.
\end{itemize}
Trip-level context embeddings are fused with sequence representations using either additive fusion or a gated fusion option.

\textbf{Key observation:} direct city prediction consistently outperformed code-based pipelines in this repository.

\subsection{Pipeline B: RQ-VAE Code Prediction}
Cities are first mapped into two-level discrete semantic codes via residual vector quantization (RQ-VAE). The sequence model then predicts the next code pair \((c_1, c_2)\), and a reverse mapping (\texttt{code\_to\_cities}) decodes candidate cities. This pipeline improves semantic compression but introduces error at both code prediction and code-to-city decoding stages.

\subsection{Pipeline C: RQKMeans Code Prediction}
This pipeline replaces learned VQ with residual k-means over Word2Vec city embeddings. Similar to RQ-VAE, it predicts next code pairs and decodes back to cities. It is computationally simpler than RQ-VAE pretraining but achieved lower peak accuracy in these experiments.

\section{Feature Engineering Summary}
Feature engineering evolved significantly and was central to gains:
\begin{itemize}[leftmargin=1.5em]
    \item \textbf{Core sequence}: city token prefixes.
    \item \textbf{Temporal}: average stay length, last stay days, month bucket.
    \item \textbf{Trip-structure}: trip length, unique-city count, repetition ratio.
    \item \textbf{Geography}: last hotel country, number of unique hotel countries, cross-border count/ratio, country streak patterns.
    \item \textbf{User context}: booker country, device class, affiliate channel.
    \item \textbf{Optional semantic side IDs}: injected RQ-derived code IDs into embedding models.
\end{itemize}

\textbf{Important data-quality correction:} context features were revised to avoid future-information leakage by computing prefix-aligned features per sample rather than full-trip statistics during training.

\section{Experimental Outcomes}
Table~\ref{tab:results} summarizes the best reported values from the project notes.

\begin{table}[h]
\centering
\caption{Best Accuracy@4 by pipeline and backbone (reported).}
\label{tab:results}
\begin{tabular}{lcc}
\toprule
\textbf{Pipeline} & \textbf{Transformer Best} & \textbf{GRU Best} \\
\midrule
Direct Embedding & 0.4871 & 0.4891 \\
RQ-VAE Codes & 0.3384 & 0.3487 \\
RQKMeans Codes & 0.3066 & 0.3291 \\
\bottomrule
\end{tabular}
\end{table}

\subsection{Practical Lessons}
\begin{itemize}[leftmargin=1.5em]
    \item \textbf{Embedding route is strongest}: direct city classification remained dominant.
    \item \textbf{Hidden-state handling matters}: correcting sequence final-state extraction improved both GRU and Transformer.
    \item \textbf{Pooling choice matters}: \texttt{last}/\texttt{mean} performed far better than \texttt{cls} in this setup.
    \item \textbf{Feature additions can help}: geographic and trip-structure features improved the embedding pipeline.
    \item \textbf{Semantic side info is nuanced}: RQ semantic IDs sometimes help slightly, but fusion strategy needs careful tuning.
\end{itemize}

\section{Engineering and Reproducibility}
The repository uses modular scripts and package organization:
\begin{itemize}[leftmargin=1.5em]
    \item \texttt{scripts/} for training entry points per pipeline.
    \item \texttt{src/features}, \texttt{src/models}, \texttt{src/training}, \texttt{src/datasets}, \texttt{src/utils}.
    \item Unified launcher support and optional direct script execution.
\end{itemize}

For reproducibility, key recommendations are:
\begin{enumerate}[leftmargin=1.5em]
    \item Freeze Python package versions in \texttt{requirements.txt}.
    \item Record seeds and hardware context for each run.
    \item Keep generated artifacts (\texttt{output/}) out of Git and version only code + configs.
\end{enumerate}

\section{Conclusion and Future Work}
This project demonstrates that, for this dataset, \textbf{direct city-token modeling with rich context features} is a strong and robust baseline, outperforming the current generative code-decoding alternatives. RQ-style routes remain promising for semantic compression and retrieval-style reasoning, but they currently trail in end-to-end Accuracy@4.

High-impact future directions include:
\begin{itemize}[leftmargin=1.5em]
    \item Validation-based early stopping and hyperparameter sweeps.
    \item Better calibrated fusion of sequence and context (e.g., stronger gated conditioning).
    \item Improved code-to-city decoding and reranking for RQ pipelines.
    \item Ensemble methods across embedding and code-based models.
\end{itemize}

\end{document}
